\chapter{Diagramas de arquitectura}
\label{ap:diagramas}

Este apéndice recoge los diagramas estructurales de la segunda fase
referenciados en el capítulo de diseño. No enumeran clases ni módulos de bajo
nivel; resumen las fronteras que hacen que el laboratorio pueda cargar modelos
generados por la primera fase, ejecutarlos, observarlos y devolver conocimiento
al backbone compartido.

\begin{figure}[p]
\centering
\resizebox{\linewidth}{!}{%
\begin{tikzpicture}[
  node distance=7mm and 10mm,
  layer/.style={
    rectangle, rounded corners=2pt, draw=black!70, line width=0.6pt,
    fill=#1, minimum height=8mm, text width=34mm,
    align=center, font=\small\sffamily
  },
  store/.style={
    cylinder, shape border rotate=90, aspect=0.25, draw=black!65,
    line width=0.55pt, fill=#1, minimum height=9mm, text width=21mm,
    align=center, font=\scriptsize\sffamily
  },
  service/.style={
    rectangle, rounded corners=2pt, draw=black!65, line width=0.55pt,
    fill=#1, minimum height=8mm, text width=26mm,
    align=center, font=\scriptsize\sffamily
  },
  note/.style={
    rectangle, rounded corners=2pt, draw=black!45, dashed, line width=0.45pt,
    fill=black!2, text width=30mm, align=center, font=\scriptsize\sffamily
  }
]
  \node[layer=blue!8] (ui) {Cliente web\\\footnotesize React, Vite, Tailwind};
  \node[layer=blue!4, below=of ui] (api) {Backend FastAPI\\\footnotesize REST + WebSocket};
  \node[layer=violet!10, below=of api] (orch) {Orchestrator\\\footnotesize sesión, tools, presupuesto};

  \node[service=green!9, below left=15mm and 28mm of orch] (architect)
       {Architect\\\footnotesize spec de entorno};
  \node[service=green!9, below=15mm of orch] (tracker)
       {Tracker\\\footnotesize eventos y trayectorias};
  \node[service=green!9, below right=15mm and 28mm of orch] (analyst)
       {Analyst\\\footnotesize patrones y contraste};
  \node[service=green!9, right=8mm of analyst] (reporter)
       {Reporter\\\footnotesize informe PDF};

  \node[layer=orange!10, below=17mm of tracker, text width=40mm] (runtime)
       {Runtime de simulación\\\footnotesize Environment 2D + modelos fase 1};

  \node[store=orange!7, left=14mm of runtime] (models) {Modelos\\aceptados};
  \node[store=red!7, right=14mm of runtime] (artifacts) {Informes\\y trazas};

  \node[store=gray!9, below left=17mm and 20mm of runtime] (pg) {PostgreSQL\\metadatos};
  \node[store=gray!9, below=17mm of runtime] (minio) {MinIO\\binarios};
  \node[store=gray!9, below right=17mm and 20mm of runtime] (qdrant) {Qdrant\\memorias};
  \node[store=gray!9, right=13mm of qdrant] (neo4j) {Neo4j\\grafo};

  \node[note, left=10mm of architect] (input) {Petición\\conversacional};
  \node[note, right=10mm of reporter] (output) {Respuesta final\\+ PDF};

  \draw[diagarrowbi] (ui) -- (api);
  \draw[diagarrow] (api) -- (orch);

  \draw[diagarrow] (input) -- (orch);
  \draw[diagarrow] (orch) -- (architect);
  \draw[diagarrow] (orch) -- (tracker);
  \draw[diagarrow] (orch) -- (analyst);
  \draw[diagarrow] (orch) -- (reporter);
  \draw[diagarrow] (reporter) -- (output);

  \draw[diagarrow] (architect) -- (runtime);
  \draw[diagarrowbi] (runtime) -- (tracker);
  \draw[diagarrow] (tracker) -- (analyst);
  \draw[diagarrow] (analyst) -- (reporter);

  \draw[diagarrow] (models) -- (runtime);
  \draw[diagarrow] (reporter) |- (artifacts);
  \draw[diagarrow] (tracker) |- (artifacts);

  \draw[diagarrow] (orch) .. controls +(15mm,-38mm) and +(0,18mm) .. (pg);
  \draw[diagarrow] (artifacts) -- (minio);
  \draw[diagarrow] (analyst) .. controls +(20mm,-24mm) and +(0,16mm) .. (neo4j);
  \draw[diagarrowbi] (analyst) -- (qdrant);
  \draw[diagarrowbi] (analyst) -- (neo4j);

  \begin{scope}[on background layer]
    \node[diaggroup, fit=(architect)(tracker)(analyst)(reporter)] {};
    \node[diaggroup, fit=(pg)(minio)(qdrant)(neo4j),
      label={[font=\scriptsize\itshape,black!60]below:Knowledge Backbone compartido}] {};
  \end{scope}
\end{tikzpicture}}
\caption{Arquitectura lógica de la segunda fase.}
\label{fig:arquitectura}

\vspace{1.6em}

\centering
\resizebox{0.66\textwidth}{!}{%
\begin{tikzpicture}[
  node distance=9mm and 12mm,
  box/.style={
    rectangle, rounded corners=2pt, draw=black!70, line width=0.6pt,
    fill=#1, minimum height=9mm, text width=34mm,
    align=center, font=\small\sffamily
  },
  db/.style={
    cylinder, shape border rotate=90, aspect=0.25, draw=black!65,
    line width=0.55pt, fill=#1, minimum height=10mm, text width=26mm,
    align=center, font=\scriptsize\sffamily
  }
]
  \node[box=green!9] (tracker) {Tracker\\\footnotesize observador puro};
  \node[box=violet!8, below=of tracker] (writer) {\texttt{Tracker}\\\texttt{MemoryWriter}\\\footnotesize consolida observaciones};
  \node[box=red!8, below=of writer] (namespace) {Memoria indexada\\\footnotesize \texttt{simulation}};

  \node[db=gray!9, below left=13mm and 24mm of namespace] (pg) {PostgreSQL\\sesión, run, IDs};
  \node[db=gray!9, below=13mm of namespace] (qdrant) {Qdrant\\denso + BM25};
  \node[db=gray!9, below right=13mm and 24mm of namespace] (minio) {MinIO\\reportes y trazas};

  \node[box=blue!8, below=24mm of qdrant] (retrieve) {Recuperación posterior\\\footnotesize Orchestrator + Analyst};
  \node[box=orange!10, below=of retrieve] (future) {Nueva sesión\\\footnotesize contexto previo};

  \draw[diagarrow] (tracker) -- node[right, font=\scriptsize] {JSON estructurado} (writer);
  \draw[diagarrow] (writer) -- node[right, font=\scriptsize] {chunking + metadatos} (namespace);
  \draw[diagarrow] (namespace) -- (pg);
  \draw[diagarrow] (namespace) -- (qdrant);
  \draw[diagarrow] (namespace) -- (minio);
  \draw[diagarrow] (qdrant) -- node[left, font=\scriptsize] {top-k / rerank} (retrieve);
  \draw[diagarrow] (pg) |- (retrieve);
  \draw[diagarrow] (minio) |- (retrieve);
  \draw[diagarrow] (retrieve) -- (future);

  \begin{scope}[on background layer]
    \node[diaggroup, fit=(pg)(qdrant)(minio),
      label={[font=\scriptsize\itshape,black!60]below:evidencia persistida por ejecución}] {};
  \end{scope}
\end{tikzpicture}}
\caption{Flujo de escritura y recuperación de observaciones de simulación.}
\label{fig:memoria-simulacion}
\end{figure}
